\documentclass[11pt]{article}
\usepackage[margin=1in]{geometry}
\usepackage{amsmath,amssymb}
\title{Minimal scalar graviton EFT through canonical dimension nine}
\author{Jarvis reproducible symbolic workbench}
\date{2026-09-05}
\begin{document}
\maketitle
Use $g_{\mu\nu}=\eta_{\mu\nu}+\kappa h_{\mu\nu}$, $[h]=1$, $[\kappa]=-1$. A vertex $\kappa^nh^n\partial^d$ has canonical operator dimension $n+d$. The result is for one minimally coupled scalar, $\xi=0$, in $d=4-2\epsilon$ and $\overline{\rm MS}$; coefficients are defined by the Euclidean determinant and continued to the mostly-minus Lorentzian convention.
\[
\Gamma_E^{(1)}=\int d^4x\sqrt g\left[C_0+C_RR+C_{R^2}R^2+C_{R_{\mu\nu}^2}R_{\mu\nu}R^{\mu\nu}+C_6\mathcal P_6+O(\partial^8)\right],
\]
\[
C_0=\frac{m^4}{64\pi^2}(L_m-\tfrac32),\quad C_R=-\frac{m^2}{192\pi^2}(L_m-1),\quad C_{R^2}=\frac{L_m}{3840\pi^2},\quad C_{R_{\mu\nu}^2}=\frac{L_m}{1920\pi^2},\quad C_6=-\frac1{32\pi^2m^2}.
\]
Here $L_m=\log(m^2/\mu^2)$. The curvature-squared result uses $E_4=R_{\mu\nu\rho\sigma}^2-4R_{\mu\nu}^2+R^2$: the minimal-scalar bulk combination is $R^2/120+R_{\mu\nu}^2/60$. The complete $\mathcal P_6$ basis, its exact coefficients, and its multiplication by $C_6$ are defined in the single canonical source \texttt{sakharov\_scalar\_graviton\_dim9.wls}.
\[
\begin{aligned}
\mathcal L_{\leq9}={}&C_0\sum_{n=0}^{9}\kappa^n[\sqrt{-g}]_{h^n}+C_R\sum_{n=2}^{7}\kappa^n[\sqrt{-g}R]_{h^n}\\
&+C_{R^2}\sum_{n=2}^{5}\kappa^n[\sqrt{-g}R^2]_{h^n}+C_{R_{\mu\nu}^2}\sum_{n=2}^{5}\kappa^n[\sqrt{-g}R_{\mu\nu}^2]_{h^n}\\
&+C_6\left(\sum_{n=2}^{3}\kappa^n[\sqrt{-g}\mathcal P_6^{(2)}]_{h^n}+\kappa^3[\sqrt{-g}\mathcal P_6^{(3)}]_{h^3}\right).
\end{aligned}
\]
Thus the retained sectors are $h^0$ through $h^9$ for the volume term, $h^2\partial^2$ through $h^7\partial^2$ for Einstein gravity, $h^2\partial^4$ through $h^5\partial^4$ for curvature squared, $h^{2,3}\partial^6$ for the quadratic part of $\mathcal P_6$, and $h^3\partial^6$ for its cubic part. The Einstein linear term is a boundary term; the volume term has a tadpole unless flat space is tuned by a counterterm.
\[
\sqrt{-g}=\exp\!\left[\frac12\sum_{r=1}^{9}\frac{(-1)^{r+1}}r\kappa^r\operatorname{tr}(\eta^{-1}h)^r\right]+O(h^{10}).
\]
The included Wolfram workbench derives every volume coefficient through $h^9$ exactly. No $B_8$ input is needed: non-total-derivative curvature-dimension-eight terms begin at $h^2\partial^8$, canonical dimension ten, while linear $\Box^3R$ terms are boundary terms.
\paragraph{Evidence.} Vassilevich, \emph{Heat kernel expansion: user's manual}, electronic PDF p.~40, eqs.~(4.26)--(4.29), at \texttt{knowledge/papers/vassilevich-heat-kernel-2003\_\_hep-th\_0306138.pdf}.
\end{document}
